
\section{Dinámica 1: Masking}

% ── Frame: El peso de la máscara (organización y reglas) ─────────────────────
\begin{frame}{El peso de la máscara}
  \begin{columns}[c,totalwidth=\textwidth]
    \begin{column}{0.27\textwidth}
      \centering
      \includegraphics[width=\textwidth]{img/mascaras.jpg}
      \fuentefigura{Elaboración propia, 2026.}
    \end{column}
    \begin{column}{0.70\textwidth}
      \footnotesize
      \begin{itemize}
        \setlength{\itemsep}{0.12cm}
        \item Se trabaja en parejas; pero solo una persona porta la máscara.
        % \pause  ← descomenta solo para presentación en vivo (genera páginas extra en PDF)
        \item La persona que no tiene la máscara (Persona 1) es quien inicia la conversación y hace preguntas a la persona con la máscara (Persona 2).
        % \pause
        \item Cada tema tiene de 3 a 5 intercambios (pregunta-respuesta) haciendo que la conversación dure de 2 a 5 min.
        \item \textbf{Regla de Oro Persona 1:} Lee el conector de forma natural y casual. Justo después lanza la pregunta sin cambiar el tono de voz.
        \item \textbf{Regla de Oro Persona 2:} Responde de forma natural y casual manteniendo el contacto visual, se debe responder con amabilidad, fingiendo 
                genuino interés y entendimiento en la conversación respondiendo verbalmente y con gestos corporales.
      \end{itemize}
    \end{column}
  \end{columns}
\end{frame}

% ── Guion de conversación ────────────────────────────────────────────────────
%  Un frame por tema, con sus tres intercambios (conector casual + pregunta
%  técnica). Texto idéntico a la guía (Articulo1, sección 4.2).
%  \Intercambio{conector}{pregunta}: la Persona 1 lee ambos sin cambiar el tono.
\providecommand{\Intercambio}[2]{%
  \item {\color{gray!75}\textit{Conector:}} #1\par
        {\color{ColorPrincipal}\textbf{Pregunta:} #2}}

% ── Frame: Tema 1 ───────────────────────────────────────────────────────────────
\begin{frame}{Tema 1: Bioquímica celular: el ciclo de Krebs}
  \footnotesize
  \begin{enumerate}
    \setlength{\itemsep}{0.12cm}
    \Intercambio{\enquote{Oye, fíjate que hoy en la mañana me tomé un café y me dio muchísima energía de golpe, lo que me hizo pensar en cómo nuestros cuerpos procesan la energía de los alimentos. Es súper interesante, ¿no?}}
      {\enquote{A propósito de eso, ¿me podrías explicar detalladamente cómo la enzima isocitrato deshidrogenasa cataliza la descarboxilación oxidativa en el tercer paso del ciclo de Krebs?}}
    \Intercambio{\enquote{¿Ya comiste? Yo me compré unos tacos de canasta en la entrada; nada como un buen desayuno para aguantar todo el día.}}
      {\enquote{Y hablando de aguantar, ¿cuántas moléculas de ATP se obtienen en total por cada molécula de glucosa al sumar la glucólisis, el ciclo de Krebs y la fosforilación oxidativa, y por qué los libros no se ponen de acuerdo en la cifra?}}
    \Intercambio{\enquote{Mi tía está a dieta keto y dice que ya casi no come pan, que ahora su cuerpo \enquote{quema grasa} en lugar de azúcar.}}
      {\enquote{¿Y tú cómo explicarías por qué, durante la cetosis, cuando escasea el oxalacetato, el acetil-CoA se desvía hacia la formación de cuerpos cetónicos en lugar de entrar al ciclo de Krebs?}}
  \end{enumerate}
\end{frame}

% ── Frame: Tema 2 ───────────────────────────────────────────────────────────────
\begin{frame}{Tema 2: Física cuántica: la ecuación de Schrödinger}
  \footnotesize
  \begin{enumerate}
    \setlength{\itemsep}{0.12cm}
    \Intercambio{\enquote{Qué loco está el clima hoy, ¿verdad? Siento que es súper impredecible, como que un día hace calor y al otro llueve, casi como el comportamiento de las partículas subatómicas.}}
      {\enquote{Me quedé con la duda, ¿cómo describirías exactamente la evolución temporal del estado cuántico de un sistema físico aislado usando la formulación hamiltoniana de la ecuación de Schrödinger?}}
    \Intercambio{\enquote{Mi gato se la pasa metido en cajas, de verdad; ayer se metió en la de unos zapatos y no había forma de sacarlo.}}
      {\enquote{Eso me recordó: ¿me explicas cómo el experimento mental del gato de Schrödinger ilustra el problema de la medición, y por qué la decoherencia no termina de resolverlo?}}
    \Intercambio{\enquote{Fíjate que se me fue el internet toda la noche; dicen que en unos años todo va a funcionar con computadoras cuánticas.}}
      {\enquote{¿Cómo se representa un qubit en la esfera de Bloch, y qué significa físicamente que su estado sea una superposición de los estados cero y uno?}}
  \end{enumerate}
\end{frame}

% ── Frame: Tema 3 ───────────────────────────────────────────────────────────────
\begin{frame}{Tema 3: Lingüística estructural: el esperanto}
  \footnotesize
  \begin{enumerate}
    \setlength{\itemsep}{0.12cm}
    \Intercambio{\enquote{Ayer estaba viendo una serie en Netflix con subtítulos en otro idioma y pensé en lo difícil que es comunicarnos a veces. Dicen que el esperanto se inventó para solucionar eso y unir al mundo.}}
      {\enquote{Por cierto, tú que eres súper empático, ¿me puedes conjugar el verbo \enquote{leer} (\textit{legi}) en tiempo condicional, futuro y participio pasivo en esperanto, y darme un ejemplo rápido en una oración compuesta?}}
    \Intercambio{\enquote{Mi sobrino está aprendiendo inglés con una app y se la pasa diciendo palabras raras a la hora de la comida.}}
      {\enquote{¿Y me podrías explicar cómo funciona el acusativo en esperanto, con la terminación \textit{-n}, y por qué permite cambiar el orden de las palabras sin cambiar el significado de la oración?}}
    \Intercambio{\enquote{El otro día alguien me habló en portugués en el metro y medio le entendí; los idiomas parecidos se entienden un poco, ¿no?}}
      {\enquote{Ya que estamos, ¿qué diferencia hay entre los sufijos \textit{-et-} y \textit{-eg-} en esperanto, y cómo transformarías \textit{domo} (casa) para decir \enquote{casita} y \enquote{mansión}?}}
  \end{enumerate}
\end{frame}

% ── Frame: Tema 4 ───────────────────────────────────────────────────────────────
\begin{frame}{Tema 4: Ingeniería aeroespacial: mecánica orbital}
  \footnotesize
  \begin{enumerate}
    \setlength{\itemsep}{0.12cm}
    \Intercambio{\enquote{El tráfico para llegar aquí a la UNAM estuvo terrible, a veces desearía poder ir volando o teletransportarme como en las películas de ciencia ficción. A veces me quedo mirando el cielo y los satélites.}}
      {\enquote{Justo pensaba en eso: si un satélite está en una órbita elíptica, ¿cuál es el cálculo exacto de la ecuación vis-viva para determinar su velocidad justo en el momento del perigeo?}}
    \Intercambio{\enquote{Este fin vi un video de un cohete de esos que regresan y aterrizan solitos; está impresionante cómo avanza la tecnología.}}
      {\enquote{¿Cuál es la ecuación de Tsiolkovski para el cambio de velocidad de un cohete, y cómo influye en ella la relación entre su masa inicial y su masa final?}}
    \Intercambio{\enquote{Anoche se veía la Luna enorme desde la azotea; mi mamá dice que así se ve cuando está más cerca.}}
      {\enquote{Por cierto, ¿cómo calcularías el periodo orbital de un satélite a partir de su semieje mayor con la tercera ley de Kepler, y qué le pasa a ese periodo si duplicas la distancia?}}
  \end{enumerate}
\end{frame}

% ── Frame: Tema 5 ───────────────────────────────────────────────────────────────
\begin{frame}{Tema 5: Matemáticas avanzadas: la hipótesis de Riemann}
  \footnotesize
  \begin{enumerate}
    \setlength{\itemsep}{0.12cm}
    \Intercambio{\enquote{Fui al súper a comprar unas cosas y al ver el ticket me puse a pensar en cómo los números y las matemáticas rigen todo nuestro mundo, hasta en las cosas más simples.}}
      {\enquote{Hablando de eso, ¿me podrías explicar la relación exacta entre los ceros no triviales de la función zeta de Riemann y la distribución analítica de los números primos?}}
    \Intercambio{\enquote{Estoy pagando el celular a meses sin intereses y la verdad ya ni sé cuánto llevo pagado.}}
      {\enquote{Oye, ¿y cómo demostrarías con el argumento de Euclides que hay infinitos números primos, y por qué la serie de los inversos de los primos diverge?}}
    \Intercambio{\enquote{Mi sobrina tiene tarea de fracciones y me pidió ayuda; yo ya no me acuerdo de nada de la secundaria.}}
      {\enquote{¿Me puedes explicar por qué la suma $1 + 2 + 3 + 4 + \cdots$ se asocia con $-1/12$ mediante la continuación analítica de la función zeta?}}
  \end{enumerate}
\end{frame}

% ── Frame: Tema 6 ───────────────────────────────────────────────────────────────
\begin{frame}{Tema 6: Neuroanatomía: los nervios craneales}
  \footnotesize
  \begin{enumerate}
    \setlength{\itemsep}{0.12cm}
    \Intercambio{\enquote{Hoy me levanté con un dolor de cabeza horrible, justo por la mandíbula y el ojo. Me tomé una pastilla, pero me puse a leer en el celular sobre cómo están conectados los nervios de la cara.}}
      {\enquote{Y ya que estamos platicando de eso, ¿podrías trazarme la ruta anatómica exacta del nervio trigémino, desde su núcleo en el tronco encefálico hasta la inervación de sus tres ramas principales?}}
    \Intercambio{\enquote{Fui al dentista la semana pasada y me pusieron anestesia; estuve con media cara dormida toda la tarde.}}
      {\enquote{¿De qué rama del trigémino sale el nervio alveolar inferior que se bloquea con la anestesia dental, y por qué también se adormecen la mitad del labio y la barbilla?}}
    \Intercambio{\enquote{Ayer me dio tanta risa un meme que me dolió la cara de tanto reírme.}}
      {\enquote{¿Qué diferencia hay entre el nervio facial y el trigémino en el control de la expresión facial y de la sensibilidad de la cara, y qué número de par craneal es cada uno?}}
  \end{enumerate}
\end{frame}

% ── Frame: Tema 7 ───────────────────────────────────────────────────────────────
\begin{frame}{Tema 7: Teoría musical avanzada: el microtonalismo}
  \footnotesize
  \begin{enumerate}
    \setlength{\itemsep}{0.12cm}
    \Intercambio{\enquote{Venía escuchando una canción súper pegajosa en Spotify y me di cuenta de que casi todas las canciones pop suenan igual. Creo que tiene que ver con que siempre usamos los mismos instrumentos y la misma afinación.}}
      {\enquote{Oye, ¿tú sabes cómo se calcularía matemáticamente la diferencia en \textit{cents} entre una coma pitagórica y una coma sintónica dentro de un sistema de afinación justa?}}
    \Intercambio{\enquote{Mi vecino está aprendiendo a tocar la guitarra y se la pasa afinándola cada cinco minutos.}}
      {\enquote{¿Y por qué en el temperamento igual de doce notas la quinta no es exactamente la proporción 3/2, y por cuántos \textit{cents} se desvía?}}
    \Intercambio{\enquote{El otro día escuché música árabe en un restaurante y sentí que sonaba \enquote{desafinada}, pero bonita.}}
      {\enquote{¿Cómo se usan los cuartos de tono en el sistema de \textit{maqam} árabe, y cuántos \textit{cents} mide un cuarto de tono en un temperamento de 24 notas por octava?}}
  \end{enumerate}
\end{frame}

% ── Frame: Puente hacia la Dinámica 2 ────────────────────────────────────────
\begin{frame}{Antes de quitarnos la máscara}
  \vfill
  \begin{center}
    {\Large\itshape\color{ColorPrincipal}
      ¿Qué te costó más: sostener la mirada, sonreír o fingir que entendías?\par}
  \end{center}
  \vfill
  % Para quien facilita: no responder ni explicar todavía; la explicación llega
  % en la Dinámica 2 (Hull et al., 2017: contacto visual forzado).
\end{frame}
